\section{Proof of Theorem~\ref{theorem: sharp}}
\label{section: proof of Theorem 1.4}

The goal of this section is to prove Theorem~\ref{theorem: sharp}.
Write $t_1$ and $\theta_{t_1}$ simply by $t$ and $\theta$.
The function $Q(x)$ in the differential
equation~\eqref{00eq: ODE on C} associated to the curvature equation
\eqref{00eq: DE on C} in the case under consideration is
\begin{equation} \label{eq: Heun Q}
Q(x)=\frac{\beta_0}{x^2}+\frac{d_0}x
+\frac{\beta_1}{(x-1)^2}+\frac{d_1}{x-1}
+\frac{\beta_t}{(x-t)^2}+\frac{d_t}{x-t},
\end{equation}
where $\beta_p=\alpha_p(\alpha_p+2)/4=\bigl(\theta_p^2-1\bigr)/4$ and $d$'s
satisfy~\eqref{e0q-32}, i.e.,
\begin{equation} \label{eq: d beta}
d_0+d_1+d_t=0, \qquad
\beta_0+\beta_1+\beta_t+d_1+td_t=\beta_\infty.
\end{equation}
Before we proceed further, we note that the set
$\{\alpha_p/2+1,-\alpha_p/2\}$ is invariant under the substitution
$\alpha_p\mapsto-2-\alpha_p=-\theta_p-1$. For simplicity of discussion later on,
we assume that~$\alpha_0$ and~$\alpha_1$ satisfy
\begin{equation} \label{eq: choice of alpha}
\alpha_0=\epsilon_0\theta_0-1, \qquad
\alpha_1=\epsilon_1\theta_1-1,
\end{equation}
where $\epsilon_0$, $\epsilon_1$ are given in the assumption of Theorem~\ref{theorem: sharp}. Thus, $\alpha_0+\alpha_1+\alpha_\infty=\epsilon_0\theta_0
+\epsilon_1\theta_1+\theta_\infty-3$ is an integer.

In view of~\eqref{eq: d beta}, we regard $d_0$ and $d_1$ as
functions of $d:=d_t$ and let $Q_{d,t}(x)$ denote the rational function
$Q(t)$ in~\eqref{eq: Heun Q}. By Theorem~\ref{prop: apparent}, there
exists a polynomial $\mathcal P(d,t)\in\Q_{\theta}[d,t]$ of degree
$\theta$ in $d$ such that the differential equation
\begin{equation} \label{eq: Heun equation}
 y''(x)=Q_{d,t}(x)y(x)
\end{equation}
is apparent at $t$ if and only if $\mathcal P(d,t)=0$, where
$\Q_{\theta}=\Q(\theta_0,\theta_1,\theta_\infty)$. The degree of
$\mathcal P(d,t)$ in~$d$ is~$\theta$, but its total degree is in
general strictly larger than $\theta$ (see~\eqref{e22q:PP}).
In our proof of Theorem~\ref{theorem: sharp}, we shall introduce
another pair $(\lambda,t)$ of unknowns in place of $(d,t)$.

Let
\begin{gather} \label{eq: lambda}
 \lambda:=td_0+t\frac{\alpha_0\alpha_1}2
 +\frac{\alpha_0\alpha_t}2
 =t(t-1)d_t+t(\beta_0+\beta_1+\beta_t-\beta_\infty)
 +t\frac{\alpha_0\alpha_1}2+\frac{\alpha_0\alpha_t}2,
\end{gather}
where the second equality follows from~\eqref{eq: d beta}. Indeed,
the parameter $\lambda$ is used as the accessary parameter of a
certain Heun equation considered in \cite{Quadrilateral}, where for
the simplicity of computation, the apparent singularity was put at
$x=0$ instead of $x=t$ in this paper, and the condition for
apparentness at $x=0$ was shown to be equivalent to the vanishing of
the characteristic polynomial of a finite Jacobi matrix (see~\cite[Proposition 2.4 and equation~(2.8)]{Quadrilateral}). For the
convenience of the reader, we recall the equivalence, stated in our
setting, as follows (see also~\mbox{\cite[Lemma B.8]{Chen-Lin-Yang}}):
\eqref{eq: Heun equation} is apparent at $t$ if and only if
 $\lambda$ is an eigenvalue of the $\theta\times\theta$ matrix
 \begin{equation} \label{eq: M}
 M=M(t):=\begin{pmatrix}
 B_1 & A_1 & & & & \\
 D_2 & B_2 & A_2 & & & \\
 & D_3 & B_3 & A_3 & & \\
 & & \ddots & \ddots & \ddots & \\
 & & & D_{\theta-1} & B_{\theta-1} & A_{\theta-1} \\
 & & & & D_{\theta} & B_{\theta} \end{pmatrix}
 \end{equation}
 (undisplayed entries are all $0$), where
 \begin{gather*}
 A_j =A_j(t):=t(t-1)j(j-\theta), \\
 B_j =B_j(t):=(2t-1)(j-1)(j-2) -(j-1)[(t-1)\alpha_0+t\alpha_1
 +(2t-1)\alpha_t]
 +t\alpha'\alpha'', \\
 D_j =D_j(t):=(j-2)(j-3)-(j-2)(\alpha_0+\alpha_1+\alpha_t)
 +\alpha'\alpha'',
 \end{gather*}
 with
\[
 \alpha':=-\frac{\alpha_0+\alpha_1+\alpha_t+\alpha_\infty}2-1,
 \qquad \alpha'':=\frac{\alpha_\infty-\alpha_0-\alpha_1-\alpha_t}2.
\]
 That is,~\eqref{eq: Heun equation} is apparent at $t$ if and only if
 $\lambda$ is an eigenvalue of $M(t)$, i.e., a zero of the polynomial
\[
 P(\lambda,t):=\det(\lambda I_\theta-M(t)).
\]
 The polynomial $P(\lambda,t)$ has the following properties.

\begin{Lemma} \label{lemma: P in lambda} \qquad
 \begin{enumerate}\itemsep=0pt
 \item[$({\rm i})$] The total degree and the degree in $\lambda$ of
 $P(\lambda,t)$ are both $\theta$.
 \item[$({\rm ii})$] The roots of $P(\lambda,0)$ are
 $(j-1)(\alpha_0+\alpha_t-j+2)$, $j=1,\dots,\theta$.
 \item[$({\rm iii})$] Let $P_\infty(\lambda):=\lim_{t\to\infty}P(\lambda
 t,t)/t^\theta$. Then the roots of $P_\infty(\lambda)$ are
 \begin{gather}
 \hat\lambda_j:=(j-1)(\alpha_\infty+\alpha_t-j+2)
 -\beta_\infty-\beta_t+\beta_0+\beta_1\nonumber\\
 \hphantom{\hat\lambda_j:=}{} +\frac{\alpha_0\alpha_1}2
 -\frac{\alpha_t\alpha_\infty}2,\qquad j=1,\dots,\theta.\label{eq: roots infinity}
 \end{gather}
 \end{enumerate}
\end{Lemma}

\begin{proof}
Statement (i) follows easily from the observations that
 $\deg_t A_j(t)=2$, ${\deg_tB_j(t)=1}$, and $D_j(t)$ are
 constant polynomials.

 The second property follows immediately from the fact that
 $A_j(0)=0$ and hence the roots of~$P(\lambda,0)$ are simply
 $B_j(0)$. The third property is proved in \cite[Theorem~B.3]{Chen-Lin-Yang}. Since the conclusions are stated in different
 ways, here we provide a sketch of proof.

 Let $u=1/x$. We check directly that $y(x)$ is a solution of~\eqref{eq: Heun equation} if and only if $\wt y(u):=uy(1/u)$
 satisfies
 \begin{equation} \label{eq: wt Heun}
 \frac{{\rm d}^2}{{\rm d}u^2}\wt y(u)=\wt Q(u)\wt y(u),
 \end{equation}
 where $\wt Q(u)=u^{-4}Q(1/u)$. Clearly, this is a Fuchsian
 differential equation with Riemann scheme
\[
 \begin{pmatrix}
 0 & 1 & 1/t & \infty \\
 -\alpha_\infty/2 & -\alpha_1/2 & -\alpha_t/2 &
 -(\alpha_0/2+1) \\
 \alpha_\infty/2+1 & \alpha_1/2+1 & \alpha_t/2+1
 & \alpha_0/2 \end{pmatrix}.
\]
 Thus,
\[
 \wt Q(u)=\frac{\beta_\infty}{u^2}+\frac{\wt d_0}u
 +\frac{\beta_1}{(u-1)^2}+\frac{\wt d_1}{u-1}
 +\frac{\beta_t}{(u-1/t)^2}+\frac{\wt d_{1/t}}{u-1/t}
\]
 for some complex numbers $\wt d_0$, $\wt d_1$, and $\wt d_{1/t}$
 satisfying
\[
 \wt d_0+\wt d_1+\wt d_{1/t}=0, \qquad
 \wt d_1+\frac{\wt d_{1/t}}t+\beta_\infty+\beta_1+\beta_t=\beta_0.
\]
In fact, computing the partial fraction decomposition of $\wt Q(u)$,
 we find that
\[
 \wt d_0=2\beta_1+2t\beta_t+d_1+t^2d_t, \qquad
 \wt d_1=-2\beta_1-d_1, \qquad
 \wt d_{1/t}=-2t\beta_t-t^2d_t.
\]
 We now apply parts (i) and (ii) to~\eqref{eq: wt Heun}. Let
\[
 \wt\lambda=\frac{\wt d_0}t+\frac{\alpha_\infty\alpha_1}{2t}
 +\frac{\alpha_\infty\alpha_t}2.
\]
 By parts (i) and (ii), there is a polynomial $\wt P\bigl(\wt\lambda,1/t\bigr)$
 such that~\eqref{eq: wt Heun} is apparent at $1/t$ if and only if
 $\wt P\bigl(\wt\lambda,1/t\bigr)=0$ and the roots of $\wt P\bigl(\wt\lambda,0\bigr)$
 are
\[
 (j-1)(\alpha_\infty+\alpha_t-j+2), \qquad
 j=1,\dots,\theta.
\]
 On the other hand, we see from~\eqref{eq: d beta} and~\eqref{eq: lambda} that
 \begin{align*}\wt d_0
 &{}=t(t-1)d_t+2t\beta_t+\beta_\infty+\beta_1-\beta_0-\beta_t\\
 &{}=\lambda+t(\beta_\infty+\beta_t-\beta_0-\beta_1)-t\frac{\alpha_0\alpha_1}{2}
 -\frac{\alpha_0\alpha_t}{2}+\beta_\infty+\beta_1-\beta_0-\beta_t, \end{align*}
 so $\wt\lambda$ and $\lambda$ are related by
 {\allowdisplaybreaks
 \begin{align*}
 \wt\lambda={}&\frac\lambda t+\beta_\infty+\beta_t-\beta_0-\beta_1
 -\frac{\alpha_0\alpha_1}2+\frac{\alpha_\infty\alpha_t}2 \\
 &{}+\frac1t(\beta_\infty+\beta_1-\beta_0-\beta_t)
 +\frac1{2t}(\alpha_1\alpha_\infty-\alpha_0\alpha_t).
 \end{align*}
 }%
 Letting $t\to\infty$, we conclude that the roots of
 $P_\infty(\lambda)$ are given by~\eqref{eq: roots infinity}.
\end{proof}

\begin{Corollary}
 \label{corollary: limiting equations}
 Let $\lambda$ and $\hat\lambda_j$ be given by~\eqref{eq: lambda} and~\eqref{eq: roots infinity}, respectively.
 Let $y_+(x;\lambda,t)$ and $y_-(x;\lambda,t)$ be solutions of~\eqref{eq: Heun equation} of the form
\begin{alignat*}{3}
 & y_+(x;\lambda,t)=x^{\alpha_0/2+1}\sum_{j=0}^\infty
 c_{+,j}(\lambda,t)x^j, \qquad && c_{+,0}(\lambda,t)=1,&
\\
 & y_-(x;\lambda,t)=x^{-\alpha_0/2}\sum_{j=0}^\infty
 c_{-,j}(\lambda,t)x^j, \qquad && c_{-,0}(\lambda,t)=1.&
\end{alignat*}
 Then as $t\to\infty$ with $\lambda/t\to\hat\lambda_j$,
 $y_\pm(x;\lambda,t)$ converge to solutions of the Fuchsian
 differential equation with Riemann scheme
 \begin{equation} \label{eq: RS infinity}
 \begin{pmatrix}
 0 & 1 & \infty \\
 -\alpha_0/2 & -\alpha_1/2 & -(\hat\alpha_\infty/2+1) \\
 \alpha_0/2+1 & \alpha_1/2+1 & \hat\alpha_\infty/2
 \end{pmatrix},
 \end{equation}
 where $\hat\alpha_\infty=\alpha_\infty+\alpha_t-2j+2$.
\end{Corollary}

\begin{proof} Let $\hat d_0$, $\hat d_1$, $\hat d_t$ be the limits of $d_0$,
 $d_1$, and $d_t$ as $t\to\infty$ and $\lambda/t\to\hat\lambda_j$.
 From~\eqref{eq: lambda}, we know that
\[
 \hat d_0=(j-1)(\alpha_\infty+\alpha_t-j+2)
 -\beta_\infty-\beta_t+\beta_0+\beta_1-\frac{\alpha_t\alpha_\infty}2
\]
 and $(t-1)d_t$ converge to finite numbers as $t\to\infty$ and
 $\lambda/t\to\hat\lambda_j$. The latter implies that
 $\hat d_t=0$. Thus, we have
\[
 \hat Q_j(x):=\lim_{t\to\infty,\,\lambda/t\to\hat\lambda_j}Q_{d,t}(x)
 =\frac{\beta_0}{x^2}+\frac{\hat d_0}x
 +\frac{\beta_1}{(x-1)^2}+\frac{\hat d_1}{x-1}.
\]
 From the relation $d_0+d_1+d_t=0$, we see that $\hat d_1=-\hat d_0$
 and
 \begin{equation*}
 \begin{split}
 \hat d_1+\beta_0+\beta_1
 &{}=-(j-1)(\alpha_\infty+\alpha_t-j+2)+\beta_\infty+\beta_t
 +\frac{\alpha_t\alpha_\infty}2 \\
 &{}=\frac14(\alpha_\infty+\alpha_t-2j+2)
 (\alpha_\infty+\alpha_t-2j+4).
 \end{split}
 \end{equation*}
 It follows that
\[
 y''(x)=\hat Q_j(x)y(x)
\]
 is a Fuchsian differential equation with Riemann scheme
 given by~\eqref{eq: RS infinity}. The convergence of~$Q_{d,t}(x)$ to~$\hat Q_j(x)$ is clearly uniform on any compact subset of ${\C\setminus\{0,1\}}$.
 Therefore, $y_\pm(x;\lambda,t)$ converge to solutions of
 $y''(x)=\hat Q_j(x)y(x)$. This completes the proof.
\end{proof}

Now for $p\in\{0,1,\infty,t\}$, let $M_p(\lambda,t)$ be the monodromy
matrices of~\eqref{eq: Heun equation} around $x=p$ with respect to the
basis $(y_+(x;\lambda,t),y_-(x;\lambda,t))^t$.
%Following the proof of Theorem~\ref{theorem: coaxial
% conditions}, we know that $M_0(\lambda,t)$ and $M_1(\lambda,t)$ are
% of the form
Recalling that the Riemann scheme of~\eqref{eq: Heun equation} is
\[
\begin{pmatrix}
 0 & 1 & t & \infty \\
 -\alpha_0/2 & -\alpha_1/2 & -\alpha_t/2 & -(\alpha_\infty/2+1) \\
 \alpha_0/2+1 & \alpha_1/2+1 & \alpha_t/2+1 & \alpha_\infty/2
\end{pmatrix},
\]
we know that
\[
M_0(\lambda,t)=\M{{\rm e}^{\pi {\rm i}\alpha_0}}00{{\rm e}^{-\pi {\rm i}\alpha_0}}, \qquad
M_t(\lambda,t)=(-1)^{\theta-1}I_2.
\]
Also, the monodromy matrices satisfy $M_0M_1M_tM_\infty=I_2$.

\begin{Lemma} We have
\[
M_1(\lambda,t)=\M{{\rm e}^{\pi {\rm i}\alpha_1}}{b(\lambda,t)}{c(\lambda,t)}
{{\rm e}^{-\pi {\rm i}\alpha_1}}
\]
for some complex numbers $b(\lambda,t)$ and $c(\lambda,t)$ satisfying
$b(\lambda,t)c(\lambda,t)=0$.
\end{Lemma}

\begin{proof} Write $M_1(\lambda,t)=\SM abcd$. Note that the
 eigenvalues of $M_1$ are ${\rm e}^{\pi {\rm i}\alpha_1}$ and ${\rm e}^{-\pi
 {\rm i}\alpha_1}$ and those of $M_\infty$ are ${\rm e}^{\pi {\rm i}\alpha_\infty}$
 and ${\rm e}^{-\pi {\rm i}\alpha_\infty}$. Together with the relation
 $M_0M_1M_tM_\infty=I_2$, these informations yield
 \begin{gather*}
 a+d={\rm e}^{\pi {\rm i}\alpha_1}+{\rm e}^{-\pi {\rm i}\alpha_1}, \\
 a{\rm e}^{\pi {\rm i}\alpha_0}+d{\rm e}^{-\pi {\rm i}\alpha_0}
 =(-1)^{\theta-1}\bigl({\rm e}^{\pi {\rm i}\alpha_\infty}+{\rm e}^{-\pi {\rm i}\alpha_\infty}\bigr).
 \end{gather*}
 Eliminating $d$, we get
\[
 \bigl({\rm e}^{\pi {\rm i}\alpha_0}-{\rm e}^{-\pi {\rm i}\alpha_0}\bigr)a
 =(-1)^{\theta-1}\bigl({\rm e}^{\pi {\rm i}\alpha_\infty}+{\rm e}^{-\pi {\rm i}\alpha_\infty}\bigr)
 -\bigl({\rm e}^{\pi {\rm i}(\alpha_1-\alpha_0)}{+{\rm e}^{\pi {\rm i}(-\alpha_1-\alpha_0)}}\bigr).
\]
 Now we have $\alpha_0+\alpha_1+\alpha_\infty=\epsilon_0\theta_0
 +\epsilon_1\theta_1+\theta_\infty-3=k-3$. Thus, the condition
 $\theta\equiv k\mod 2$ implies that
\[
 {\rm e}^{\pi {\rm i}(-\alpha_1-\alpha_0)}={\rm e}^{\pi {\rm i}(3-k+\alpha_\infty)}
 =(-1)^{\theta-1}{\rm e}^{\pi {\rm i}\alpha_\infty}
\]
 and hence $(-1)^{\theta-1}{\rm e}^{-\pi {\rm i}\alpha_\infty}
 ={\rm e}^{\pi {\rm i}(\alpha_0+\alpha_1)}$.
 From these, we see that $a={\rm e}^{\pi {\rm i}\alpha_1}$. It follows that
 $d={\rm e}^{-\pi {\rm i}\alpha_1}$ and $bc=0$. This proves the lemma.
\end{proof}

Let $P(\lambda,t)=P_1(\lambda,t)\cdots P_m(\lambda,t)$ be
the factorization of $P(\lambda,t)$ into a product of irreducible
polynomials over $\C$. It is easy to see that for each $j$, we have
either $b(\lambda,t)\equiv 0$ or $c(\lambda,t)\equiv 0$ on $\mathcal
A_j:=\bigl\{(\lambda,t)\in\C^2\colon P_j(\lambda,t)=0\bigr\}$. Thus, if we set
\[
 P_b(\lambda,t):=\prod_{b(\lambda,t)\equiv 0\text{ on }\mathcal A_j}
 P_j(\lambda,t), \qquad
 P_c(\lambda,t):=\prod_{c(\lambda,t)\equiv 0\text{ on }\mathcal A_j}
 P_j(\lambda,t),
\]
then the cardinality of the set $A$ is equal to the number of
intersections of the two curves $P_b(\lambda,t)=0$ and
$P_c(\lambda,t)=0$ in the affine plane $\C^2$, with the multiplicity
of a point in $A$ defined to be that of the corresponding intersection
point of $P_b$ and $P_c$. Furthermore, by part~(iii) of Lemma~\ref{lemma: P in lambda}, the two curves $P_b(\lambda,t)=0$ and
$P_c(\lambda,t)=0$ do not intersect at infinity. Therefore, by
Bezout's theorem, the number of intersections of the two curves in the
affine plane is $(\deg P_b)(\deg P_c)$. Hence, to prove Theorem~\ref{theorem: sharp}, we only need to determine the degrees of~$P_b$
and~$P_c$.

\begin{Proposition} \label{prop: degree of Pb}
 Let $k:=\theta_\infty+\epsilon_0\theta_0+\epsilon_1\theta_1$ be
 given as in the statement of Theorem~$\ref{theorem: sharp}$.
 If~$\theta\le|k|$, then
\[
 \begin{cases}
 \deg P_b=0\text{ and }\deg P_c=\theta, &\text{if }k>0, \\
 \deg P_b=\theta\text{ and }\deg P_c=0, &\text{if }k<0.
 \end{cases}
\]
 If $\theta>|k|$, then $\deg P_b=(\theta-k)/2$ and $\deg
 P_c=(\theta+k)/2$.
\end{Proposition}

To prove Proposition~\ref{prop: degree of Pb}, we recall that, by
Lemma~\ref{lemma: P in lambda}\,(i), the total degree and the degree in
$\lambda$ of $P(\lambda,t)$ are equal. It is easy to see that every
factor of $P(\lambda,t)$ has the same property. It follows that the
degree of $P_b(\lambda,t)$ is equal to the degree of the polynomial
$P_{b,\infty}(\lambda):=\lim_{t\to\infty}P_b(\lambda t,t)/t^{\deg
 P_b}$, which in turn is equal to the number of $\hat\lambda_j$ that
are roots of $P_{b,\infty}(\lambda)$, where $\hat\lambda_j$ are given
by~\eqref{eq: roots infinity}. In other words, to determine the degree
of $P_b(\lambda,t)$, we shall count how many $\hat\lambda_j$ such that
$\lim_{t\to\infty,\lambda/t\to\hat\lambda_j}c(\lambda,t)\neq 0$ there
are.

\begin{Lemma}\label{lem-111} Let $\hat\lambda_j$ be given by~\eqref{eq: roots
 infinity} and
\[
 \hat b=\lim_{t\to\infty,\lambda/t\to\hat\lambda_j}b(r,t), \qquad
 \hat c=\lim_{t\to\infty,\lambda/t\to\hat\lambda_j}c(r,t).
\]
 Then
\[
 \begin{cases}
 \hat b=0,~\hat c\neq 0, & \text{if }j>(k+\theta)/2, \\
 \hat b\neq0,~\hat c=0, &\text{if }j\le(k+\theta)/2.
 \end{cases}
\]
\end{Lemma}

\begin{proof} By Corollary~\ref{corollary: limiting equations},
 $y_\pm(x;\lambda,t)$ converge to solutions $\hat y_\pm(x)$ of the
 Fuchsian differential equation with Riemann scheme~\eqref{eq: RS
 infinity}, where $\hat y_\pm(x)$ are local solutions near $0$ of
 the form $\hat y_+(x)=x^{\alpha_0/2+1}(1+O(x))$ and $\hat
 y_-(x)=x^{-\alpha_0/2}(1+O(x))$. Then
 $w_{0,\pm}(x):=x^{\alpha_0/2}(x-1)^{\alpha_1/2}\hat
 y_\pm(x)$ are solutions of the hypergeometric differential equation
 with Riemann scheme
\[
 \begin{pmatrix} 0 & 1 & \infty \\
 0 & 0 & \alpha \\
 1-\gamma & \gamma-\alpha-\beta & \beta\end{pmatrix},
\]
 where
\[
 \gamma=-\alpha_0, \qquad
 \alpha=-\left(\frac{\alpha_0+\alpha_1+\hat\alpha_\infty}2+1\right),
 \qquad
 \beta=\frac{\hat\alpha_\infty-\alpha_0-\alpha_1}2.
\]
 Let $w_{1,\pm}(x)$ be the local solutions near $x=1$ of the same
 hypergeometric differential equation of the form
 $w_{1,+}(x)=(x-1)^{\gamma-\alpha-\beta}(1+O(x-1))$ and
 $w_{1,-}(x)=1+O(x-1)$. By~\cite[Chapter~2, Theorem~4.7.1]{Yoshida}, the connection matrix $P$ such that
 $(w_{0,+}(x),w_{0,-}(x))^t=P{(w_{1,+}(x),w_{1,-}(x))^t}$ is
\[
 P=\begin{pmatrix}
 \displaystyle
 \frac{\Gamma(2-\gamma)\Gamma(\gamma-\alpha-\beta)}
 {\Gamma(1+\alpha-\gamma)\Gamma(1+\beta-\gamma)} &
 \displaystyle
 \frac{\Gamma(2-\gamma)\Gamma(\gamma-\alpha-\beta)}
 {\Gamma(1-\alpha)\Gamma(1-\beta)} \\[10pt]
 \displaystyle
 \frac{\Gamma(\gamma)\Gamma(\gamma-\alpha-\beta)}
 {\Gamma(\alpha)\Gamma(\beta)} &
 \displaystyle\frac{\Gamma(\gamma)\Gamma(\gamma-\alpha-\beta)}
 {\Gamma(\gamma-\alpha)\Gamma(\gamma-\beta)} \end{pmatrix}.
\]
 Then the monodromy matrix of the hypergeometric differential
 equation around $x=1$ with respect to the basis
 $(w_{0,+}(x),w_{0,-}(x))^t$ is
\[
 P\M{{\rm e}^{2\pi {\rm i}(\gamma-\alpha-\beta)}}001P^{-1}.
\]
 Now according to our choice of $\alpha_0$ and $\alpha_1$ in~\eqref{eq: choice of alpha}, we have
 \begin{align*}
 \alpha&{}=-\left(\frac{\alpha_0+\alpha_1+\hat\alpha_\infty}2
 +1\right)
=-\frac{(\epsilon_0\theta_0-1)+(\epsilon_1\theta_1-1)
 +(\theta_\infty-1+\theta-1-2j+2)}2-1 \\
 &{}=j-\frac{k+\theta}2,
 \end{align*}
 which is an integer by the assumption $\theta\equiv k\mod 2$.
 Thus, $P$ is upper-triangular when $j\le(k+\theta)/2$ and is
 lower-triangular when $j>(k+\theta)/2$. This implies that $\hat
 b\neq0$ and $\hat c=0$ if $j\le(k+\theta)/2$, and $\hat b=0$ and
 $\hat c\neq 0$ if $j>(k+\theta)/2$.
\end{proof}

\begin{proof}[Proof of Proposition~\ref{prop: degree of Pb} and Theorem~\ref{theorem: sharp}] Since $j=1,2,\dots, \theta$, Proposition~\ref{prop: degree of Pb} follows from Lemma~\ref{lem-111}. Consequently, Theorem~\ref{theorem: sharp} holds.
\end{proof}
